% Generated draft: requires mathematical and source review.
\documentclass[11pt]{article}
\usepackage[margin=25mm]{geometry}
\usepackage{fontspec}
\setmainfont{Latin Modern Roman}
\usepackage{amsmath,amssymb,mathtools,unicode-math}
\setmathfont{Latin Modern Math}
\usepackage{enumitem,needspace,xurl,hyperref}
\setlength{\parindent}{0pt}
\setlength{\parskip}{5pt}
\setlength{\emergencystretch}{3em}
\newcommand{\N}{\mathbb{N}}
\newcommand{\Z}{\mathbb{Z}}
\newcommand{\Q}{\mathbb{Q}}
\newcommand{\R}{\mathbb{R}}
\newcommand{\C}{\mathbb{C}}
\newcommand{\F}{\mathbb{F}}
\newcommand{\A}{\mathbb{A}}
\newcommand{\PP}{\mathbb{P}}
\newcommand{\E}{\mathbb{E}}
\newcommand{\eps}{\varepsilon}
\newcommand{\dd}{\,\mathrm{d}}
\begin{document}
\section*{UnsolvedMath ID: 20001939}
\textbf{Source problem code:} catalogue entry
\subsection{Definitions and mathematical statement}
Two smooth pseudo-Riemannian metrics \(g,\bar g\) on a manifold are projectively equivalent if they have the same unparameterized geodesics. Equivalently, their Levi-Civita connections satisfy
\[
 \bar\nabla_XY=\nabla_XY+\alpha(X)Y+\alpha(Y)X
\]
for a one-form \(\alpha\). They are affinely equivalent if \(\bar\nabla=\nabla\), and proportional if \(\bar g=cg\) for a nonzero constant \(c\).

Question 12.0.6 asks whether there exist two nonproportional, projectively equivalent, smooth metrics of indefinite signature on \(S^3\). In dimension three an indefinite nondegenerate metric has signature \((2,1)\) or \((1,2)\); the question concerns metrics on the entire closed sphere, not merely on a coordinate domain. No completeness assumption is added: pseudo-Riemannian geodesic completeness does not follow merely from compactness.

The supplied research aid reports the following affine obstruction. If \(M\) is a closed connected Lorentzian three-manifold with finite fundamental group, every pseudo-Riemannian metric \(h\) on \(M\) sharing the Levi-Civita connection of \(g\) has the form \(h=cg\). Thus a hypothetical pair answering the question on \(S^3\) would necessarily be non-affine.

For clarity, a standard normalization of the compatibility tensor in dimension three is
\[
 L=\left|\frac{\det\bar g}{\det g}\right|^{1/4}\bar g^{-1}g.
\]
It is \(g\)-self-adjoint and satisfies the linear compatibility (Sinjukov) equation
\[
 \nabla_kL^i{}_j=\delta^i_k\lambda_j+g^{i\ell}\lambda_\ell g_{jk},
 \qquad \lambda=\tfrac12\dd(\operatorname{tr}L).
\]
Define the degree of mobility \(D(g)\) as the dimension over \(\R\) of the space of smooth \(g\)-self-adjoint solutions of this equation on the whole manifold, including possibly noninvertible solutions. The aid reports, using this equation and unspecified established global results, that any hypothetical nonproportional projectively equivalent indefinite pair on \(S^3\) must satisfy
\[
 \bar\nabla\ne\nabla,\qquad
 \dd(\operatorname{tr}L)\not\equiv0,\qquad
 D(g)=2,\qquad
 \operatorname{Spec}(L_x)\subset\R\quad\text{for every }x\in S^3.
\]
Real spectrum here does not assert diagonalizability or distinct eigenvalues. These are the aid's reduction claims; the supplied record gives no precise references or hypotheses of the global theorems used in obtaining them.

A separate reported endpoint-sign argument excludes the standard diagonal orthogonal Hopf-coordinate Lorentz ansatz: it cannot extend smoothly as a Lorentzian metric across both singular circles of the Hopf-coordinate description of \(S^3\). The record supplies neither its coefficient formula nor its exact ansatz hypotheses. This restricted exclusion is therefore retained as reported, without claiming to exclude general metrics with mixed coordinate terms or arbitrary projectively equivalent pairs. Neither obstruction supplies an example or a general nonexistence theorem for Question 12.0.6.
\subsection{Short English statement}
The question asks for two globally defined indefinite metrics on the three-sphere that share unparameterized geodesics without being constant multiples. The aid rules out the affine case and reports further necessary conditions on any remaining example, plus an obstruction to one diagonal coordinate construction. The general existence question remains unresolved in this supplied attempt.
\subsection{Sources}
\begin{itemize}
\item Local UnsolvedMath registry, permanent ID 20001939, exact source code catalogue entry; canonical \texttt{aim-geometry-notes.json}, \texttt{notes[276]}, tag \texttt{question}; AIM Workshop Problem Lists, Geometry, Geodesics, source item 12.0, Question 12.0.6. Created and updated 14 August 2026.
\item Keith Burns and Vladimir Matveev, \emph{Open problems and questions about geodesics}, last question; \url{https://aimath.org/pastworkshops/geodesicsproblems.pdf}, as identified by the supplied statement verification. The linked PDF was not accessed. No precise external references for the aid's global projective-equivalence results or affine-rigidity theorem are supplied.
\item Supplied attempt 1: result status \texttt{partial\_result}, problem status at run \texttt{open}, rigor label \texttt{proved}, novelty confidence \texttt{low}; classification \texttt{PARTIAL-PROGRESS}, registry status \texttt{partially\_solved}. Candidate contribution: affine obstruction, mobility-two real-spectrum reduction, and diagonal Hopf-coordinate exclusion. These are dated source assessments, not independently verified theorems.
\item The appended acknowledgements thank workshop participants Victor Bangert, Charles Boubel, Nancy Hingston, Gerhard Knieper, Yiming Long, Gabriel Paternain, Mark Pollicott, Regina Rotman, Stéphane Sabourau, and Benjamin Schmidt, and also Juan-Carlos Alvarez Paiva, Ivan Babenko, Florent Balacheff, Helga Baum, Michael Bialy, Dima Burago, Maciej Dunajski, Misha Gromov, Larry Guth, Sergei Ivanov, Bruce Kleiner, Alan Reid, and Stefan Suhr. They mention three visits by KB to Friedrich-Schiller-Universität Jena supported by DFG (SPP 1154 and GK 1523). These are acknowledgements, not hypotheses or individual attribution of Question 12.0.6.
\item Ulam AI / UnsolvedMath Contributors; CC BY 4.0. Machine-generated research aid, not peer reviewed; independently verify before citation.
\end{itemize}
\end{document}
